% GENERATED FILE. Edit canonical lessons, then run npm run generate:books.
% canonical-source-hash: fnv1a64:8042f32874a24154
% canonical-lessons: FR-C447-pressing, FR-C447-laverie, FR-C447-detacher, FR-C447-teindre, FR-C447-retoucher, FR-R447-revision-440-447

\chapter{Au pressing}
\label{ch:fr-au-pressing}
\hlchaptermodality{447}

\begin{chapteropening}
\textbf{By the end of this chapter:} \emph{I can get clothes cleaned and altered: the dry cleaner, the launderette, removing stains, dyeing and altering.}

Complete the last lesson of chapter 447: I can get clothes cleaned and altered: the dry cleaner, the launderette, removing stains, dyeing and altering.
\end{chapteropening}

\section[le pressing]{le pressing — the dry cleaner’s}

\label{lesson:FR-C447-pressing}



\begin{quote}
Before the new one: say the French for the locksmith, then the French for to sew back on.
\end{quote}

\subsection*{You'll want to know: le pressing}
\textbf{le pressing} — \textquotedblleft{}the dry cleaner’s\textquotedblright{}.

\begin{quote}
Il dépose sa veste au pressing le lundi.
\end{quote}

\begin{quote}
\emph{He drops his jacket off at the dry cleaner's on Mondays.}
\end{quote}

\emph{Le nettoyage à sec} — dry cleaning.

\begin{cousinweb}
\emph{pressing} was borrowed from English \emph{pressing} 'ironing', a sense English does not use for a shop. EtymDB does not list it.
\end{cousinweb}

\begin{grammarlens}[title={What you've built}]
One more word for telling, explaining and describing at B1.
\end{grammarlens}

\subsection*{Your turn}
\begin{itemize}
\raggedright
  \item \emph{Say it:} \emph{le pressing}
  \item \emph{Say it:} \emph{le pressing}, once more
  \item \emph{Say it:} a sentence of your own with \emph{le pressing}
  \item \emph{Recall:} say the French for the locksmith, then the French for to sew back on, then say \emph{le pressing} again
\end{itemize}

\begin{tcolorbox}[breakable,colback=teal!4,colframe=teal!35!black,title={Before you move on}]
What does le pressing mean? (\textquotedblleft{}The dry cleaner’s\textquotedblright{}.) Say it once more.
\end{tcolorbox}
\section[la laverie]{la laverie — the launderette}

\label{lesson:FR-C447-laverie}



\begin{quote}
Before the new one: say the French for to sew back on, then the French for the dry cleaner’s.
\end{quote}

\subsection*{You'll want to know: la laverie}
\textbf{la laverie} — \textquotedblleft{}the launderette\textquotedblright{}.

\begin{quote}
Les étudiants lavent leur linge à la laverie.
\end{quote}

\begin{quote}
\emph{The students wash their clothes at the launderette.}
\end{quote}

\emph{Une laverie automatique}; in Québec, \emph{une buanderie}.

\begin{cousinweb}
\emph{laverie} is built on \emph{laver} 'to wash', which is inherited from Latin \emph{lavāre} (Lewis \& Short). EtymDB does not list it.
\end{cousinweb}

\begin{grammarlens}[title={What you've built}]
One more word for telling, explaining and describing at B1.
\end{grammarlens}

\subsection*{Your turn}
\begin{itemize}
\raggedright
  \item \emph{Say it:} \emph{la laverie}
  \item \emph{Say it:} \emph{la laverie}, once more
  \item \emph{Say it:} a sentence of your own with \emph{la laverie}
  \item \emph{Recall:} say the French for to sew back on, then the French for the dry cleaner’s, then say \emph{la laverie} again
\end{itemize}

\begin{tcolorbox}[breakable,colback=teal!4,colframe=teal!35!black,title={Before you move on}]
What does la laverie mean? (\textquotedblleft{}The launderette\textquotedblright{}.) Say it once more.
\end{tcolorbox}
\section[détacher]{détacher — to remove a stain from; to untie}

\label{lesson:FR-C447-detacher}



\begin{quote}
Before the new one: say the French for the dry cleaner’s, then the French for the launderette.
\end{quote}

\subsection*{You'll want to know: détacher}
\textbf{détacher} — \textquotedblleft{}to remove a stain from; to untie\textquotedblright{}.

\begin{quote}
Le pressing a bien détaché ma chemise blanche.
\end{quote}

\begin{quote}
\emph{The dry cleaner got the stain out of my white shirt.}
\end{quote}

Two senses: from \emph{tache} 'stain', and from \emph{attacher} — \emph{détacher sa ceinture}, to undo one's belt.

\begin{cousinweb}
\emph{détacher} in the sense 'to untie' comes from Old French \emph{destachier} (EtymDB), the opposite of \emph{attacher}; in the sense 'to remove a stain' it is built on \emph{tache} 'stain'.
\end{cousinweb}

\begin{grammarlens}[title={What you've built}]
One more word for telling, explaining and describing at B1.
\end{grammarlens}

\subsection*{Your turn}
\begin{itemize}
\raggedright
  \item \emph{Say it:} \emph{détacher}
  \item \emph{Say it:} \emph{détacher}, once more
  \item \emph{Say it:} a sentence of your own with \emph{détacher}
  \item \emph{Recall:} say the French for the dry cleaner’s, then the French for the launderette, then say \emph{détacher} again
\end{itemize}

\begin{tcolorbox}[breakable,colback=teal!4,colframe=teal!35!black,title={Before you move on}]
What does détacher mean? (\textquotedblleft{}To remove a stain from; to untie\textquotedblright{}.) Say it once more.
\end{tcolorbox}
\section[teindre]{teindre — to dye}

\label{lesson:FR-C447-teindre}



\begin{quote}
Before the new one: say the French for the launderette, then the French for to remove a stain from.
\end{quote}

\subsection*{You'll want to know: teindre}
\textbf{teindre} — \textquotedblleft{}to dye\textquotedblright{}.

\begin{quote}
Elle a teint son vieux pull en bleu.
\end{quote}

\begin{quote}
\emph{She dyed her old jumper blue.}
\end{quote}

Irregular: \emph{je teins, j'ai teint}. \emph{La teinturerie} — another word for the dry cleaner's.

\begin{cousinweb}
\emph{teindre} is inherited through Old French \emph{taindre} from Latin \emph{tingere} 'to dye, to wet' (EtymDB).
\end{cousinweb}

\begin{grammarlens}[title={What you've built}]
One more word for telling, explaining and describing at B1.
\end{grammarlens}

\subsection*{Your turn}
\begin{itemize}
\raggedright
  \item \emph{Say it:} \emph{teindre}
  \item \emph{Say it:} \emph{teindre}, once more
  \item \emph{Say it:} a sentence of your own with \emph{teindre}
  \item \emph{Recall:} say the French for the launderette, then the French for to remove a stain from, then say \emph{teindre} again
\end{itemize}

\begin{tcolorbox}[breakable,colback=teal!4,colframe=teal!35!black,title={Before you move on}]
What does teindre mean? (\textquotedblleft{}To dye\textquotedblright{}.) Say it once more.
\end{tcolorbox}
\section[retoucher]{retoucher — to alter (a garment); to touch up}

\label{lesson:FR-C447-retoucher}



\begin{quote}
Before the new one: say the French for to remove a stain from, then the French for to dye.
\end{quote}

\subsection*{You'll want to know: retoucher}
\textbf{retoucher} — \textquotedblleft{}to alter (a garment); to touch up\textquotedblright{}.

\begin{quote}
Le tailleur a retouché le bas du pantalon.
\end{quote}

\begin{quote}
\emph{The tailor altered the bottom of the trousers.}
\end{quote}

\emph{Une retouche} — an alteration; also on photos.

\begin{cousinweb}
\emph{retoucher} is built from \emph{re-} 'again' and \emph{toucher} 'to touch' (EtymDB lists the verb), which is usually traced to a Vulgar Latin verb that imitates a knock.
\end{cousinweb}

\begin{grammarlens}[title={What you've built}]
One more word for telling, explaining and describing at B1.
\end{grammarlens}

\subsection*{Your turn}
\begin{itemize}
\raggedright
  \item \emph{Say it:} \emph{retoucher}
  \item \emph{Say it:} \emph{retoucher}, once more
  \item \emph{Say it:} a sentence of your own with \emph{retoucher}
  \item \emph{Recall:} say the French for to remove a stain from, then the French for to dye, then say \emph{retoucher} again
\end{itemize}

\begin{tcolorbox}[breakable,colback=teal!4,colframe=teal!35!black,title={Before you move on}]
What does retoucher mean? (\textquotedblleft{}To alter (a garment); to touch up\textquotedblright{}.) Say it once more.
\end{tcolorbox}
\section[(dialogue)]{Retour sur les chapitres 440–447}

\label{lesson:FR-R447-revision-440-447}



\begin{quote}
Say the French for the old suburb and the district.
\end{quote}

\subsection*{Your turn}
Say these aloud:

\begin{itemize}
\raggedright
  \item the old suburb, the district, the prefecture, the police station, the barracks
  \item the bench, the street lamp, the lamp post, the kiosk, the cobblestone
  \item the façade, the shop window, the shop sign, the roof gutter, the covered market
  \item the sewer, the gutter, to clear away, to salt, to clear of snow
  \item the pork butcher, the fishmonger’s, the caterer, the grocer, the greengrocer
  \item the hardware shop, the household goods shop, the stationer’s, the haberdasher’s, the wine merchant
  \item the cobbler, to resole, to polish, the locksmith, to sew back on
  \item the dry cleaner’s, the launderette, to remove a stain from, to dye, to alter
\end{itemize}

\begin{tcolorbox}[breakable,colback=teal!4,colframe=teal!35!black,title={Before you move on}]
The old suburb? (\textbf{le faubourg}.) The district? (\textbf{l'arrondissement}.) The prefecture? (\textbf{la préfecture}.) The police station? (\textbf{le commissariat}.) The barracks? (\textbf{la caserne}.) The bench? (\textbf{le banc}.) The street lamp? (\textbf{le lampadaire}.) The lamp post? (\textbf{le réverbère}.) The kiosk? (\textbf{le kiosque}.) The cobblestone? (\textbf{le pavé}.) The façade? (\textbf{la façade}.) The shop window? (\textbf{la vitrine}.) The shop sign? (\textbf{l'enseigne}.) The roof gutter? (\textbf{la gouttière}.) The covered market? (\textbf{la halle}.) The sewer? (\textbf{l'égout}.) The gutter? (\textbf{le caniveau}.) To clear away? (\textbf{déblayer}.) To salt? (\textbf{saler}.) To clear of snow? (\textbf{déneiger}.) The pork butcher? (\textbf{la charcuterie}.) The fishmonger’s? (\textbf{la poissonnerie}.) The caterer? (\textbf{le traiteur}.) The grocer? (\textbf{l'épicier}.) The greengrocer? (\textbf{le primeur}.) The hardware shop? (\textbf{la quincaillerie}.) The household goods shop? (\textbf{la droguerie}.) The stationer’s? (\textbf{la papeterie}.) The haberdasher’s? (\textbf{la mercerie}.) The wine merchant? (\textbf{le caviste}.) The cobbler? (\textbf{le cordonnier}.) To resole? (\textbf{ressemeler}.) To polish? (\textbf{cirer}.) The locksmith? (\textbf{le serrurier}.) To sew back on? (\textbf{recoudre}.) The dry cleaner’s? (\textbf{le pressing}.) The launderette? (\textbf{la laverie}.) To remove a stain from? (\textbf{détacher}.) To dye? (\textbf{teindre}.) To alter? (\textbf{retoucher}.)
\end{tcolorbox}
